\documentclass{article}

\usepackage{amssymb,amsmath}

\begin{document}

General model assumptions:
\begin{itemize}
\item the fluid density and viscosity are constant in time and space (implies incompressibility);
\item the fluid flow is laminar;
\item the column is operated under constant conditions (e.g., temperature, flow rate);
\item the process is isothermal (i.e., there are no thermal effects);
\item diffusion does not depend on the concentration of the components, viscosity of the fluid, or pressure;
\item the fluid in the particles is stagnant (i.e., there is no convective flow);
\item the porous particles are spherical, rigid, and do not move;
\item the particles are homogeneous and of uniform porosity;
\item all components access the same particle volume;
\item the partial molar volumes are the same in mobile and solid phase;
\item the binding model parameters do not depend on pressure and are constant along the column;
\item the solvent is not adsorbed;
\end{itemize}


Specific model assumptions:
\begin{itemize}
\item the particles form a continuum inside the column (i.e., there is interstitial and particle volume at every point in the column);
\item the column is radially symmetric and homogeneous (i.e., concentration profiles and parameters only depend on the axial position);
\item the interstitial liquid phase concentration is spatially constant on the outer particle surface;
\item there is no surface diffusion (i.e. adsorbed molecules are spatially constant);
\end{itemize}


Linear binding model assumptions:
\begin{itemize}
\item each component binds independently, that is, there are no competitive effects;
\item the capacity of the resin is unbounded;
\end{itemize}


Michaelis Menten reaction model assumptions:
\begin{itemize}
\item the reaction rates follow Michaelis-Menten enzyme kinetics, characterized by saturation at high substrate concentrations;
\item the kinetic expressions are interpreted in the steady-state (kinetic) sense;
\item this model applies to reactions occurring within a single phase (liquid or solid) only;
\end{itemize}


\section*{General Rate Model with Linear binding with Michaelis Menten type reactions}
Consider a cylindrical column of length $L > 0$ packed with spherical particles, and observed over a time interval $(0, T^{\mathrm{end}})$.
In the interstitial volume, mass transfer is governed by the following convection-diffusion-reaction equations in $(0, T^\mathrm{end})\times (0, L)$ and for all components $i\in\{1, \dots, N^{\mathrm{c}} \}$
\begin{align}
\varepsilon^{\mathrm{c}} \frac{\partial c^{\mathrm{b}}_i}{\partial t} = - u \frac{\partial \left( \varepsilon^{\mathrm{c}} c^{\mathrm{b}}_i \right)}{\partial z} + \frac{\partial}{\partial z} \left( \varepsilon^{\mathrm{c}} D^{\mathrm{ax}}_{i} \frac{\partial c^{\mathrm{b}}_i}{\partial z} \right)- \left(1 - \varepsilon^{\mathrm{c}} \right) \frac{3}{R^{\mathrm{p}}} k^{\mathrm{f}}_{i} \left(c^{\mathrm{b}}_i - \left. c^{\mathrm{p}}_{i} \right|_{r = R^{\mathrm{p}}} \right) + f^{\mathrm{react},\mathrm{b}}_{i}\left( \vec{c}^{\mathrm{b}} \right),
\end{align}
with boundary conditions

\begin{alignat}{2}
u c_{\mathrm{in},i} &= \left.\left( u c^{\mathrm{b}}_i - D^{\mathrm{ax}}_{i} \frac{\partial c^{\mathrm{b}}_i}{\partial z}\right)\right|_{z=0} & &\qquad\text{on }(0, T^{\mathrm{end}}),\\
               0 &= - D^{\mathrm{ax}}_{i} \left. \frac{\partial c^{\mathrm{b}}_i}{\partial z} \right|_{z=L} & &\qquad\text{on }(0, T^{\mathrm{end}}).
\end{alignat}
Here, $t\in (0, T^{\mathrm{end}})$ is the time coordinate, $z\in (0, L)$ is the axial cylinder coordinate, $R^\mathrm{p}> 0$ is the particle radius, $c^{\mathrm{b}}_i\colon (0, T^\mathrm{end})\times (0, L) \mapsto \mathbb{R}$ is the bulk liquid concentration, $\varepsilon^{\mathrm{c}}\in (0, 1)$ is the column porosity, $u> 0$ is the interstitial velocity, $D^\mathrm{ax}_i\geq 0$ is the axial dispersion coefficient, $k^\mathrm{f}_{i}\geq 0$ is the film diffusion coefficient, $N^{\mathrm{react},\mathrm{b}}$ is the number of bulk reactions, $S^{\mathrm{b}}_{i,r}$ is the stoichiometric matrix (bulk), $v^{\mathrm{max},\mathrm{b}}_{r}$ is the maximum reaction rate (bulk), $K^{\mathrm{M},\mathrm{b}}_{m,r}$ is the Michaelis constant (bulk), and $N^{\mathrm{sub},\mathrm{b}}_{r}$ is the number of substrates per reaction (bulk).
In the particles, mass transfer is governed by diffusion-reaction equations in $(0, T^\mathrm{end}) \times (0, L)\times (0, R^{\mathrm{p}})$ and for all components

\begin{align}
\frac{\partial c^{\mathrm{p}}_{i}}{\partial t} &= \frac{1}{r^2} \frac{\partial }{\partial r} \left( r^2 D_{i}^{\mathrm{p}} \frac{\partial c^{\mathrm{p}}_{i}}{\partial r} \right) - \frac{1 - \varepsilon^{\mathrm{p}}}{\varepsilon^{\mathrm{p}}}\left(k^{\mathrm{a}}_{i} c^{\mathrm{p}}_{i} - k^{\mathrm{d}}_{i} c^{\mathrm{s}}_{i} \right), \\
\frac{\partial c^{\mathrm{s}}_{i}}{\partial t} &= k^{\mathrm{a}}_{i} c^{\mathrm{p}}_{i} - k^{\mathrm{d}}_{i} c^{\mathrm{s}}_{i} ,
\end{align}

with boundary conditions
\begin{align}
- \left. \left( D^{\mathrm{p}}_{i} \frac{\partial c^{\mathrm{p}}_{i}}{\partial r} \right) \right|_{r=0}
&= 0, \\
\varepsilon^{\mathrm{p}} \left. \left( D^{\mathrm{p}}_{i} \frac{\partial c^{\mathrm{p}}_{i}}{\partial r} \right)\right|_{r = R^{\mathrm{p}}_{}}
               &= k^{\mathrm{f}}_{i} \left. \left( c^{\mathrm{b}}_i - c^{\mathrm{p}}_{i} \right|_{r = R^{\mathrm{p}}_{}} \right).\end{align}
Here, $r$ is the radial particle coordinate, $c^{\mathrm{p}}_{i}\colon (0, T^\mathrm{end}) \times (0, L)\times (0, R^{\mathrm{p}}) \mapsto \mathbb{R}$ is the particle liquid concentration, $c^{\mathrm{s}}_{i}\colon (0, T^\mathrm{end}) \times (0, L)\times (0, R^{\mathrm{p}}) \mapsto \mathbb{R}$ is the particle solid concentration, $\varepsilon^{\mathrm{p}}\in (0, 1)$ is the particle porosity, $D^\mathrm{p}_{i}> 0$ is the particle diffusion coefficient, $k^{\mathrm{a}}_{i}$ is the adsorption rate constant, and $k^{\mathrm{d}}_{i}$ is the desorption rate constant.
The bulk liquid phase reaction term is defined by Michaelis Menten type reactions (original formulation Michaelis and Menten (1913), steady-state interpretation Briggs and Haldane (1925), modern notation and parameter definitions Cornish-Bowden (2012)).

\begin{align}
f^{\mathrm{react},\mathrm{b}}_{i}\left( \vec{c}^{\mathrm{b}} \right) = \sum_{r=0}^{N^{\mathrm{react},\mathrm{b}}-1} S^{\mathrm{b}}_{i,r} \,\nu^{\mathrm{b}}_{r}
\end{align}

where

\begin{align}
\nu^{\mathrm{b}}_{r} = v^{\mathrm{max},\mathrm{b}}_{r} \prod_{m=1}^{N^{\mathrm{sub},\mathrm{b}}_{r}} \frac{c^{\mathrm{b}}_{m,r}}{K^{\mathrm{M},\mathrm{b}}_{m,r} + c^{\mathrm{b}}_{m,r}}.
\end{align}

Consistent initial values for all solution variables (concentrations) are defined at $t = 0$.
\begin{thebibliography}{99}
\bibitem{Michaelis1913} Michaelis, Leonor and Menten, Maud L. (1913). Die Kinetik der Invertinwirkung. Biochemische Zeitschrift.
\bibitem{Briggs1925} Briggs, G. E. and Haldane, J. B. S. (1925). A Note on the Kinetics of Enzyme Action. Biochemical Journal.
\bibitem{CornishBowden2012} Cornish-Bowden, Athel (2012). Fundamentals of Enzyme Kinetics.
\end{thebibliography}
\end{document}